% 注：main.tex 中译注宏为 \znote / \znoteplain（任务书写作 \zhnote），此处按 main.tex 实际定义使用。
\chapter{模曲线作为椭圆曲线的模空间}

% ===== PDF p.055 =====
\noindent{\small\emph{（本讲由 Ben Heuer 撰写）}}

\textbf{上次}：
\begin{itemize}
\item 朗兰兹互反 $\mathrm{LL}\colon\cA_{\mathrm{mod}}\xrightarrow{\ \simeq\ }\mathcal{G}_{\mathrm{mod}}$ 的构造的初步草图：
\[
\cA_{\mathrm{mod}}:=\Bigl\{\pi\subseteq\bigoplus_{k\ge1}S_k\ \Bigm|\ \pi\ \text{是不可约的}\ \GL_2(\Af)\text{-子表示}\Bigr\},
\]
\[
\mathcal{G}_{\mathrm{mod}}:=\Bigl\{\sigma\colon G_\Q\to\GL_2(\Ql)\ \Bigm|\ \text{在有限集}\ S\ \text{之外非分歧、奇、在}\ \ell\ \text{上方 de Rham}\Bigr\}.
\]
\item 构造的关键点：对充分小的紧开子群 $K\subseteq\GL_2(\Af)$，复流形
\[
\GL_2(\Q)\backslash\bigl(\GL_2(\Af)/K\times\HH^{\pm}\bigr)
\]
(1) 在 $\Spec(\C)$ 上是代数的，(2) 有 $\Spec(\Q)$ 上的典范模型 $X_K$。由此得到 $\ell$-进平展上同调 $H^1_{\mathrm{\acute{e}t}}(X_{K,\overline{\Q}},\Q_\ell)$ 上的 Galois 作用。\znoteplain{「$K$ 充分小」是用来消去椭圆曲线的非平凡自同构的：它使商 $K\backslash\HH^{\pm}$ 成为光滑复曲线；再由 Baily--Borel 定理，这个商其实是拟射影代数簇，因而带有 $\Spec(\C)$ 上的代数结构。}
\end{itemize}

\textbf{今天}：
\begin{itemize}
\item 解释代数模型 $X_K$ 为何存在。
\item 把此前定义的概念用代数几何的语言重新解释：模形式、紧化、$q$-展开、adélic Hecke 作用，等等。
\end{itemize}

\medskip
\noindent\textbf{椭圆曲线的回顾}

\noindent 参考文献：[Silverman: The Arithmetic...]
\begin{itemize}
\item 设 $K$ 为任意域。椭圆曲线有如下相互等价的定义。
\end{itemize}

\noindent\textbf{定义 7.1}\phantomsection\label{def:7.1}\quad 域 $K$ 上的椭圆曲线等价地是下列对象之一：
\begin{itemize}
\item 亏格为 $1$ 的连通光滑射影曲线 $E|K$，连带一个选定的点 $O\in E(K)$；
\item $1$ 维连通光滑射影代数群 $E|K$；
\item 由 Weierstraß 方程
\[
E\colon y^2=x^3+ax+b,\qquad a,b\in K
\]
定义的非奇异平面三次曲线（若 $\operatorname{char}K\notin\{2,3\}$ 即为此式；否则需要更多项，参见 [Del75]）。
\end{itemize}

\begin{itemize}
\item \textbf{事实}：上述群概形自动是交换的（！）。
\item \textbf{事实}：非奇异性可以写成判别式的条件：
\[
E\ \text{非奇异}\iff\Delta(a,b)=-16(4a^3+27b^2)\ne0\qquad(\text{当}\ \operatorname{char}K\notin\{2,3\}).
\]
\item \textbf{事实}：给定 $E$，Weierstraß 方程并不唯一。但 $j$-不变量
\[
j(E):=j(a,b)=1728\,\frac{4a^3}{4a^3+27b^2}
\]
与 Weierstraß 方程的选取无关。
\item \textbf{事实}：若 $K$ 代数闭，则
\[
E\cong E'\iff j(E)=j(E').
\]
\item 更一般地，可以把 $\Spec(K)$ 换成任意基概形 $S$。
\end{itemize}

\noindent\textbf{定义 7.2}\phantomsection\label{def:7.2}\quad $S$ 上的椭圆曲线是指一个固有光滑曲线 $E\to S$，其几何纤维连通且亏格为 $1$，并带有一个点 $0\colon S\to E$。

% ===== PDF p.056 =====
\begin{itemize}
\item 可以把它看作以 $S$ 为参数的椭圆曲线族。
\item 它同样自动带有群结构：
\end{itemize}

\noindent\textbf{定理}（Abel）\quad 给椭圆曲线 $E\to S$ 赋予交换 $S$-群概形的结构、使得 $0$ 为单位截面，这样的方式唯一。

\begin{itemize}
\item 在 $S$ 上 Zariski 局部地，可以用一个 Weierstraß 方程描述 $E$。
\end{itemize}

\noindent\textbf{定义 7.3}\phantomsection\label{def:7.3}\quad 椭圆曲线之间的态射 $E\to E'$（在 $S$ 上）是指 $S$-群概形的同态。\textbf{同源\idx{同源}}是指 fppf-满射的同态，且其核有限平坦。

\noindent\textbf{例 7.4}\phantomsection\label{ex:7.4}\quad 对任意 $N\in\Z$，群概形结构给出「乘以 $N$」的态射，记作 $[N]\colon E\to E$。
\begin{itemize}
\item \textbf{事实}：$[N]$ 是秩为 $N^2$ 的有限平坦态射。若 $N$ 在 $S$ 上可逆，则它是有限平展的。
\item 特别地，$[N]$ 是同源。
\end{itemize}

\noindent\textbf{定义 7.5}\phantomsection\label{def:7.5}\quad 记 $E[N]:=\ker[N]\subseteq E$，即 $N$-挠子群。这是有限平坦的闭子群概形；当 $N$ 在 $S$ 上可逆时，它是有限平展的。
\begin{itemize}
\item 更一般地，对任意指数为 $N$ 的有限平坦子群概形 $D\subseteq E[N]$，存在唯一的椭圆曲线 $E/D$ 以及核为 $D$ 的同源 $\phi\colon E\to E/D$。
\item $[N]$ 诱导出对偶同源 $\phi^\vee\colon E/D\to E$，其核为 $E[N]/D$。
\end{itemize}

\medskip
\noindent\textbf{$K=\C$ 上的椭圆曲线}

\noindent 参考文献：[Diamond--Shurman: A first course...]
\begin{itemize}
\item 设 $E|\C$ 为椭圆曲线。则 $E(\C)$ 具有紧复 Lie 群的结构。作为这样的对象，它可以被典范地单值化：
\item 记 $\Lie E$ 为 $0\in E(\C)$ 处的切空间。则 $\Lie E\cong\C$，且有复 Lie 群的同构
\[
E(\C)=\Lie E/H_1(E,\Z)\cong\C/\Lambda,\qquad \Lambda\cong\Z^2 .
\]
\item 对任意 $N\in\Z$，有
\[
E(\C)[N]=\bigl(\tfrac1N\Lambda\bigr)/\Lambda\subseteq\C/\Lambda .
\]
\item \textbf{注}：$E(\C)[N]=\bigl(\tfrac1N\Z/\Z\bigr)^2=(\Z/N)^2$ 的秩为 $N^2$，与上述相符。
\item \textbf{要点}：$E(\C)$ 是 $1$ 维复环面：即 $\C$ 对一个格 $\Z^2\cong\Lambda\subseteq\C$（满足 $\Lambda\otimes_\Z\R=\C$）的商。
\item 复环面的态射 $:=$ 复 Lie 群的态射。
\end{itemize}

\noindent\textbf{定理}\quad 有范畴等价
\[
\{\C\ \text{上的椭圆曲线}\}\longrightarrow\{1\ \text{维复环面}\}.
\]
\begin{itemize}
\item 本质满性：设 $\Lambda\subseteq\C$ 是如上之格。定义 Weierstraß $\wp$ 函数
\[
\wp(z,\Lambda)=\frac1{z^2}+\sum_{w\in\Lambda,\ w\ne0}\Bigl(\frac1{(z-w)^2}-\frac1{w^2}\Bigr).
\]
\znoteplain{底本提取文本此处把括号内写作 $1/z^2-1/(z-w)^2$；求和时该式在 $|w|\to\infty$ 发散；Weierstraß $\wp$ 函数的标准定义是 $1/(z-w)^2-1/w^2$，减去主项后和式才对 $|w|\to\infty$ 收敛，译文从后者。}
\item 于是 $\wp(z)=\wp(z,\Lambda)$ 在 $\C\backslash\Lambda$ 上全纯，其复导数为
\[
\wp'(z)=-2\sum_{w\in\Lambda}\frac1{(z-w)^3}.
\]
\end{itemize}

% ===== PDF p.057 =====
\begin{itemize}
\item 它显然以 $\Lambda$ 为周期，$\rightsquigarrow$ $\wp$ 与 $\wp'$ 是 $\C/\Lambda$ 上的亚纯函数。
\end{itemize}

\noindent\textbf{定理}\quad $\wp$ 与 $\wp'$ 由一个 Weierstraß 方程相联系：
\[
E_\Lambda\colon \wp'^2=4(\wp^3+g_2(\Lambda)\wp+g_3(\Lambda)),
\]
其中 $g_2(\Lambda),g_3(\Lambda)\in\C$ 有显式表达式。于是得到 Lie 群的同构
\[
(\wp,\wp')\colon \C/\Lambda\xrightarrow{\ \sim\ }E_\Lambda(\C).
\]
\begin{itemize}
\item 把 $\C/\Lambda\mapsto E_\Lambda$ 送过去，定义了从 $\{1\ \text{维复环面}\}$ 到 $\{\C\ \text{上的椭圆曲线}\}$ 的拟逆。
\end{itemize}

\medskip
\noindent\textbf{椭圆曲线的复模空间}
\begin{itemize}
\item 现在我们可以把复椭圆曲线按同构加以参数化。
\item 为此只需确定何时两个格 $\Lambda_1,\Lambda_2$ 满足 $\C/\Lambda_1\cong\C/\Lambda_2$。
\item \textbf{事实}：复环面之间的任何同态都是乘以某个 $a\in\C$。事实上，它可以提升为 Lie 代数上的全纯同态。
\item 于是这发生当且仅当 $\Lambda_1=a\Lambda_2$ 对某个 $a\in\C^\times$ 成立。
\item 现在可以如下单值化：写 $\Lambda=w_1\Z\oplus w_2\Z$。
\item 条件 $\Lambda\otimes_\Z\R=\C$ 蕴含 $w_1/w_2\in\C\backslash\R=\HH^\pm$。
\item 从而对每个复环面 $E$，存在 $\tau\in\HH^\pm$ 使得
\[
E\cong E_\tau:=\C/\Lambda_\tau,\qquad \Lambda_\tau=\Z\oplus\tau\Z .
\]
\item 最后，对 $\tau,\tau'\in\HH^\pm$，有
\[
\Z+\tau\Z=\Z+\tau'\Z\iff\exists\,\gamma=\begin{pmatrix}a&b\\c&d\end{pmatrix}\in\GL_2(\Z):\ \gamma\tau=\frac{a\tau+b}{c\tau+d}=\tau'.
\]
\item 这表明：$\tau\mapsto E_\tau$ 定义了
\[
\GL_2(\Z)\backslash\HH^\pm\cong\{\text{复环面}\}/\!\sim\ \cong\{\C\ \text{上的椭圆曲线}\}/\!\sim .
\]
\item 这给我们早先的空间
\[
\SL_2(\Z)\backslash\HH=\GL_2(\Z)\backslash\HH^\pm=\GL_2(\Q)\backslash\bigl(\GL_2(\Af)/K\times\HH^\pm\bigr),\qquad K=\GL_2(\widehat{\Z})
\]
以「模空间」的解释（目前还是弱意义上的）：
\[
\SL_2(\Z)\backslash\HH=\{\C\ \text{上的椭圆曲线}\}/\!\sim .
\]
\item 由于 $\C$ 代数闭，$j$-不变量给出双射（事实上是双全纯）
\[
j\colon \SL_2(\Z)\backslash\HH\xrightarrow{\ \sim\ }\C .
\]
\item 通过覆叠映射 $\HH^\pm\to\GL_2(\Z)\backslash\HH^\pm$，由上述得到：
\[
\HH^\pm=\{\text{带有序基}\ \alpha\colon\Z^2\xrightarrow{\sim}\Lambda\ \text{的复环面}\}/\!\sim\ =\{\C\ \text{上带}\ \alpha\colon\Z^2\xrightarrow{\sim}H_1(E,\Z)\ \text{的椭圆曲线}\ E\}/\!\sim .
\]
\item 其他级的模解释呢？回忆对 $N\in\N$，我们曾有级 $K=K(N)$，
\[
X_K:=\GL_2(\Q)\backslash\bigl(\GL_2(\Af)/K\times\HH^\pm\bigr)=\coprod\Gamma(N)\backslash\HH^\pm,
\]
其中
\[
\Gamma(N)=\Bigl\{\gamma\in\GL_2(\Z)\ \Bigm|\ \gamma\equiv\begin{pmatrix}1&0\\0&1\end{pmatrix}\bmod N\Bigr\}.
\]
\end{itemize}

% ===== PDF p.058 =====
\begin{itemize}
\item 这些正是保持子群 $E_\Lambda[N]=\frac1N\Lambda/\Lambda\subseteq\C/\Lambda$ 的 $\Lambda$ 的自同构。
\item 于是有更一般的陈述
\[
\begin{aligned}
\Gamma(N)\backslash\HH^\pm&=\Bigl\{\text{带有序基}\ \alpha\colon(\Z/N)^2\xrightarrow{\sim}\tfrac1N\Lambda/\Lambda\\
&\text{的复环面}\Bigr\}/\!\sim\\
&=\Bigl\{\C\ \text{上带}\ \alpha\colon(\Z/N)^2\xrightarrow{\sim}E(\C)[N]\ \text{的椭圆曲线}\ E\Bigr\}/\!\sim .
\end{aligned}
\]
\item 在 $\GL_2(\Z)$ 与 $\Gamma(N)$ 之间，我们曾有群
\[
\begin{gathered}
\Gamma_1(N)=\Bigl\{\gamma\in\GL_2(\Z)\ \bigm|\ \gamma\equiv\left(\begin{smallmatrix}*&*\\[1.5pt]0&1\end{smallmatrix}\right)\bmod N\Bigr\},\\
\Gamma_0(N)=\Bigl\{\gamma\in\GL_2(\Z)\ \bigm|\ \gamma\equiv\left(\begin{smallmatrix}*&*\\[1.5pt]0&*\end{smallmatrix}\right)\bmod N\Bigr\}.
\end{gathered}
\]
\item 与它们相联系的模曲线有模解释：
\[
\Gamma_1(N)\backslash\HH^\pm=\{\C\ \text{上的椭圆曲线}\ E\ \text{连带一个恰为}\ N\ \text{阶的点}\ Q\in E(\C)[N]\}/\!\sim,
\]
\[
\Gamma_0(N)\backslash\HH^\pm=\{\C\ \text{上的椭圆曲线}\ E\ \text{连带一个秩为}\ N\ \text{的循环子群}\ C\subseteq E(\C)[N]\}/\!\sim .
\]
\end{itemize}

\medskip
\noindent\textbf{引导例：Legendre 族}
\begin{itemize}
\item 考虑变量 $\lambda$ 的复流形 $U=\Pone\backslash\{0,1,\infty\}$。
\item 在 $U$ 上，我们有一个复椭圆曲线族：
\item Legendre 族由 Weierstraß 方程
\[
E_\lambda\colon Y^2=X(X-1)(X-\lambda)
\]
从 $\mathbb{P}^2_U\to U$ 中割出。
\end{itemize}

\noindent\textbf{命题 7.6}\phantomsection\label{prop:7.6}
\begin{enumerate}
\item[(1)] 有 $E_\lambda[2]=\{O,(0,0),(1,0),(\lambda,0)\}$。特别地，有自然同构 $E_\lambda[2]=(\Z/2\Z)^2$。
\item[(2)] 对 $\C$ 上任意椭圆曲线 $E$ 连同平凡化 $\alpha\colon(\Z/2\Z)^2\xrightarrow{\sim}E[2]$，存在唯一的点 $x\in U$ 使得 $E=(E_\lambda)_x$。
\end{enumerate}

\begin{itemize}
\item \textbf{要点}：$E_\lambda\to U$ 在某种意义下是带 $\Gamma(2)$-级的万有椭圆曲线。
\item 这一切都是代数的！可以把 $E_\lambda\to U$ 理解为 $\C$ 上概形的态射。
\item 更好的是：这些概形已经定义在 $\Q$ 上。
\item 可以把这一论证细化，证明
\[
\Gamma(2)\backslash\HH^\pm=:X_{\Gamma(2)}(\C)=\Pone\backslash\{0,1,\infty\}.
\]
\item \textbf{注}（NB）：特别地，紧化为 $X^*_{\Gamma(2)}=\Pone$。
\item 这表明 $S_2(\Gamma_0(2))=0$，这是我们在 Fermat 大定理中用到的。\footnote{（原注 11）我们隐含地用了 $\omega^{\otimes2}\cong\Omega^1$。}
\end{itemize}

\noindent\textbf{思路}：从 $\C$ 过渡到一般 $\Q$-概形 $S$ 上椭圆曲线的模，以得到 $\Gamma(N)\backslash\HH^\pm$ 在 $\Q$ 上的代数模型。

\begin{itemize}
\item \textbf{问题}：在 $\Q$ 上的概形范畴上，是否存在一个概形表示函子
\[
P_{\SL_2(\Z)}\colon S\mapsto\{S\ \text{上的椭圆曲线}\}/\!\sim\ ?
\]
这将是一个 $\Q$ 上的概形 $X$，其 $\C$-点为 $\SL_2(\Z)\backslash\HH$。\znoteplain{原文把「不带级结构」的模函子记作 $P_{\SL_2(\Z)}$（底本提取文本作 PSL2(Z)，下标丢失）；这里 $P_{-}$ 是模函子的记号，并不是群 $\mathrm{PSL}_2(\Z)$。}
\end{itemize}

% ===== PDF p.059 =====
\begin{itemize}
\item \textbf{回答}：不存在。看出这一点的一个简便途径是扭（twist）现象：
\item 可以有在 $\Q$ 上不同构、而在某个二次扩域 $K$ 上同构的两条椭圆曲线 $E\not\cong E'$。特别地，$j(E)=j(E')$。
\item 显然，$X$ 的两个 $\Q$-点相同当且仅当它们在 $K$ 上相同。所以这样的 $X$ 不可能存在！
\item 上面的 Legendre 族有同样的问题：它并不表示
\[
S\mapsto\{S\ \text{上的椭圆曲线}\ E\ \text{连同}\ \alpha\colon(\Z/2\Z)^2\xrightarrow{\sim}E[2]\}/\!\sim,
\]
因为它没有把二次扭考虑进去。
\item \textbf{事实}：域 $K$ 上椭圆曲线 $E$ 的扭对应于 1-余圈 $G_{\overline{K}|K}\to\Aut(E)$。二次扭来自对合 $[-1]\colon E\to E\in\Aut(E)$。
\item \textbf{口号}：非平凡的自同构对可表性是不利的。\znoteplain{原文脚注（12）：更确切地说，出问题的是自同构群的非平坦性（由椭圆曲线的例外同构造成）。例如，在许多情形下，参数化线丛同构类的 Picard 函子是可表的；但在诸如固有、几何整纤维的假设下，线丛的自同构群都是 $\Gm$，从而是平坦的。}
\item \textbf{要点}：为了得到可表的函子，需要额外的数据。
\end{itemize}

\medskip
\noindent\textbf{椭圆曲线的模问题}

\noindent [Katz--Mazur: Arithmetic moduli...]
\begin{itemize}
\item 考虑 $\Z[\frac1N]$ 上概形上的模函子：
\[
P_{\Gamma_0(N)}\colon S\mapsto\{(E|S,\ C\subseteq E[N]\ \text{为秩}\ N\ \text{的循环子群概形})\}/\!\sim,
\]
\[
P_{\Gamma_1(N)}\colon S\mapsto\{(E|S,\ Q\in E[N](S)\ \text{为恰}\ N\ \text{阶的点})\}/\!\sim,
\]
\[
P_{\Gamma(N)}\colon S\mapsto\{(E|S,\ \alpha\colon(\Z/N\Z)^2\xrightarrow{\sim}E[N]\ \text{为群概形的同构})\}/\!\sim .
\]
\item 也可以构造「混合的模问题」，例如
\[
P_{\Gamma_1(N)\cap\Gamma_0(p)}:=P_{\Gamma_1(N)}\times_{P_{\SL_2(\Z)}}P_{\Gamma_0(p)}\colon S\mapsto\{(E,Q,C)\}/\!\sim .
\]
\item 关键点：设 $(E,C)\in P_{\Gamma_0(N)}(S)$。自同构 $[-1]\colon E\to E$ 把 $C$ 送到 $C$。特别地，$[-1]\in\Aut(E,C)$。
\item 相反，设 $(E,Q)\in P_{\Gamma_1(N)}(S)$，则 $[-1]\colon E\to E$ 把 $Q\in E[N]$ 送到 $-Q\in E[N]$；当 $N\ge3$ 时二者不同。特别地，$[-1]\notin\Aut(E,Q)$。\znoteplain{原文脚注（13）：而且没有例外同构。}
\item $P_{\Gamma(N)}(S)$ 的情形类似。这解释了 Legendre 族的问题：$N=2$ 太小，对付不了自同构 $[-1]$！我们需要 $N\ge3$。
\item \textbf{要点}：当 $N\ge3$ 时，模问题 $P_{\Gamma(N)}$ 与 $P_{\Gamma_1(N)}$ 是「刚性的」（$:=$ 没有非平凡自同构），而 $P_{\Gamma_0(N)}$ 与 $P_{\Gamma(2)}$ 则不然。
\end{itemize}

\noindent\textbf{定理}\quad 设 $N\ge3$，$p$ 为任意素数。
\begin{enumerate}
\item[(1)] 模问题 $P_{\Gamma(N)}$ 与 $P_{\Gamma_1(N)}$ 分别可表为 $\Z[\frac1N]$ 上的光滑仿射曲线 $X_{\Gamma(N)}$ 与 $X_{\Gamma_1(N)}$。
\item[(2)] 对任意 $n\in\N$，模问题 $P_{\Gamma(N)\cap\Gamma_0(p^n)}$ 可表为 $\Z[\frac1N]$ 上的平坦仿射曲线 $X_{\Gamma_1(N)\cap\Gamma_0(p^n)}$，且该曲线在 $\Z[\frac1{pN}]$ 上光滑。\znoteplain{原文此处把模问题 $P_{\Gamma(N)\cap\Gamma_0(p^n)}$ 的可表曲线记作 $X_{\Gamma_1(N)\cap\Gamma_0(p^n)}$，下标与模问题不一致；按上下文二者应同为 $\Gamma(N)\cap\Gamma_0(p^n)$。此处照录原文记号。}
\end{enumerate}

% ===== PDF p.060 =====
\begin{itemize}
\item \textbf{注}：为了表示 $P_{\SL_2(\Z)}$，可以从概形过渡到更大的叠范畴。得到「椭圆曲线的模叠」$\cM_{1,1}$。
\item 在各级结构 $\Gamma(N)\subseteq\Gamma'\subseteq\Gamma$ 之间，有遗忘态射 $X_{\Gamma'}\to X_\Gamma$。它们在 $\Z[\frac1N]$ 上都是有限平展的。
\item $\GL_2(\Z/N\Z)$ 从右边作用在 $P_{\Gamma(N)}$ 上，通过预复合 $\alpha\mapsto\alpha\circ\gamma$。
\item 等价地，这是通过用 $\gamma^\vee=\gamma^{-1}\det\gamma$ 预复合得到的左作用。
\item 这诱导出 $X_{\Gamma(N)}$ 上的 $\GL_2(\Z/N\Z)$-作用。
\item 类似地，$X_{\Gamma_1(N)}$ 上有自然的 $(\Z/N\Z)^\times$-作用。
\item 于是可以定义模曲线
\[
X_{\Gamma_0(N)}:=X_{\Gamma_1(N)}/(\Z/N\Z)^\times .
\]
\item 它不表示 $P_{\Gamma_0(N)}$，但「尽可能地接近」。
\item 特别地，对代数闭域 $K$，有 $X_{\Gamma_0(N)}(K)=P_{\Gamma_0(N)}(K)$。
\end{itemize}

\medskip
\noindent\textbf{紧化}
\begin{itemize}
\item 由 Weierstraß 方程给出的 $j$-不变量定义了有限平坦函数
\[
j\colon X\to\A^1 .
\]
\item 把 $X$ 经 $j\colon X\to\A^1\hookrightarrow\Pone$ 看成 $\Pone$ 的开子概形，再取它在函数域 $K(X)$ 中的正规化，可以定义
\[
j\colon X^*\to\Pone,
\]
它仍是有限平坦态射。把 $X^*$ 想成 $X$ 添上一个有限除子。
\item \textbf{事实}：$X^*$ 光滑且固有。
\item \textbf{理由}：可以用 $\Z[\frac1N][\![q]\!]$ 上的 Tate 曲线看出。
\item $X^*(\C)$ 就是前面提到的 $\Gamma\backslash\HH^\pm$ 的光滑紧化 $X^*_\C$。
\item \textbf{注}：$X^*$ 是「广义椭圆曲线」的模空间，参见 [DR73]。
\end{itemize}

\medskip
\noindent\textbf{几何模形式\idx{几何模形式}}

\noindent 参考文献：[Katz]，[Loeffler: lectures notes]
\begin{itemize}
\item 固定 $N\ge3$。令 $X=X_{\Gamma_1(N)}$。
\item 目标：把模形式几何地重新解释为 $X$ 上某些层的截面：
\end{itemize}

\noindent\textbf{定义 7.7}\phantomsection\label{def:7.7}\quad 令 $\omega:=e^*\Omega^1_{E|X}$，其中 $e\colon X\to E$ 是单位截面。
\begin{itemize}
\item 由于 $E\to X$ 是 $1$ 维光滑态射，这是一个线丛。
\item \textbf{事实}：它唯一地延拓为 $X^*$ 上的线丛 $\omega$。\textbf{理由}：万有椭圆曲线延拓为群概形 $E^*\to X^*$，取 $e^*\Omega^1_{E^*|X^*}$。
\end{itemize}

\noindent\textbf{命题 7.8}\phantomsection\label{prop:7.8}\quad $\omega^{\otimes k}$ 的解析化自然同构于 $X^*(\C)$ 上权 $k$ 的模形式层 $\omega_k$。

\noindent\textbf{证明略述}．回忆：$\HH^\pm$ 是带有 $\alpha\colon\Z^2\xrightarrow{\sim}H_1(E,\Z)$ 的 $E|\C$ 的模空间。沿 $\alpha(1,0)$ 积分定义了同构 $e^*\Omega^1_{E|\C}\xrightarrow{\sim}\C$。于是得到 $\omega$ 在 $\HH^\pm$ 上的典范平凡化。验证 $\GL_2(\Z)$-作用与之相符即可。\hfill$\square$

\begin{itemize}
\item 特别地，层 $\omega$ 已经有 $\Z[\frac1N]$ 上的代数模型。
\end{itemize}

% ===== PDF p.061 =====
\noindent\textbf{定义 7.9}\phantomsection\label{def:7.9}\quad 对任意 $\Z[\frac1N]$-代数 $R$，可以定义
\[
M_k(\Gamma_1(N),R)=\Gamma\bigl(X^*_{\Gamma_1(N)}\times_{\Spec(\Z[1/N])}\Spec(R),\ \omega^{\otimes k}\bigr),
\]
即系数在 $R$ 中的权 $k$ 模形式构成的 $R$-模。
\begin{itemize}
\item \textbf{事实}：它是有限自由的 $\Z[\frac1N]$-模。
\item \textbf{事实}：利用 Tate 曲线，可以得到 $R[[q]]$ 中的 $q$-展开。$\rightsquigarrow$ 由此可以定义尖点形式。
\item \textbf{事实}：若 $R,S$ 是平坦的 $\Z[\frac1N]$-代数（例如任何 $\Q$-代数），则
\[
M_k(\Gamma_1(N),R)\otimes_R S=M_k(\Gamma_1(N),S).
\]
\item \textbf{要点}：模形式显然是代数几何的对象！
\end{itemize}
